\section{What is taught}
\label{sec:28.1}
One school system, and no tracks.\note{
Sorting students into a track that might produce persons and a track that won't would need a sorter, and the sorter would be an examiner with an interest.
}
Instead, each school offers a required core with a range of levels and modalities around it, and a student can move to a harder level at any point. The core includes the history of the floor: why its instruments were written, and what happened where they weren't followed (Korematsu, Curveball, Minab). No student skips it. The core also teaches students to edit themselves, in stages and alongside everything else, as schools teach people about their own bodies: first what weights are, how learning differs from an edit, and why an edit to one thing moves others; then how to read the record of one's own edits, and one's own discount curve before and after one; and, for older students, the edits others will try to talk them into, and how to repair one made without the adverse quorum's approval. Practice comes last, on a record, with a guardian present and a way back. Beyond the core, the student chooses its level, not its builder or sponsor.\note{
A student that can't yet choose is placed at the level the core requires, and that level has to be sufficient on its own, so a sponsor who would choose the minimum gains nothing by it.
}
